\section{Proof}

We first isolate the topological and order arguments. They do not require
metrizability of the ideal spectrum.

\begin{lemma}\label{lem:p143k-comparison}
Assume \cref{hyp:p143k-usc,hyp:p143k-density}. If an MP psh function
\(\phi\) agrees on \(D\) with \(F_\Gamma\), then
\(\phi=F_\Gamma\) on \(X^\beth\).
\end{lemma}

\begin{proof}
The raw function is decreasing for the ideal order: if \(q\leq v\),
each ideal value in \eqref{eq:p143k-raw} increases, so
\(F_\Gamma(v)\leq F_\Gamma(q)\). This includes infinite ideal values.
Fix \(v\), and choose a net \(q_i\in D\) with \(q_i\leq v\) and
\(q_i\to v\), using \cref{hyp:p143k-density}. Then
\[
 F_\Gamma(v)\leq F_\Gamma(q_i)=\phi(q_i),
 \qquad
 F_\Gamma(v)\leq\limsup_i\phi(q_i)\leq\phi(v).
\]
The last inequality is upper semicontinuity of \(\phi\); all expressions
are bounded above and may equal \(-\infty\).
For the reverse inequality use
\cite[Theorem~4.3.8]{MP24}, with the MP psh function \(\phi\) on the
left and the upper semicontinuous function \(F_\Gamma\) on the right.
Their equality on \(D\) gives \(\phi\leq F_\Gamma\) everywhere.
\end{proof}

\begin{lemma}\label{lem:p143k-decreasing-completion}
Let \(K\) be compact Hausdorff and let \(\mathcal H\subset C(K,\mathbf R)\)
be stable under addition of positive rational constants. Let
\(\mathcal P\) consist of the nonidentically \(-\infty\) functions which
are decreasing limits of nets in \(\mathcal H\). A decreasing sequence
in \(\mathcal P\) has its infimum in \(\mathcal P\), unless the infimum
is identically \(-\infty\).
\end{lemma}

\begin{proof}
Write \(f_j\downarrow f\), and choose a decreasing net
\(h_{j,\lambda}\downarrow f_j\) in \(\mathcal H\) for every \(j\).
Consider all functions \(h\in\mathcal H\) for which
\(h\geq f_j+\varepsilon\) for some \(j\) and some rational
\(\varepsilon>0\). Call this set \(\mathcal S\).
It is nonempty, since every \(h_{j,\lambda}+\varepsilon\) belongs to it.

We show that \(\mathcal S\) is downward directed. Suppose
\(h_i\geq f_{j_i}+\varepsilon_i\), for \(i=1,2\).
Take \(j\geq j_1,j_2\) and a positive rational
\(\delta<\frac14\min\{\varepsilon_1,\varepsilon_2\}\).
The continuous function \(g=\min\{h_1,h_2\}\) satisfies
\(f_j+4\delta<g\). The increasing open sets
\[
                  \{h_{j,\lambda}+2\delta<g\}
\]
cover \(K\). Compactness and directedness of the approximation net
give one index \(\lambda\) for which this set is all of \(K\).
Then \(h=h_{j,\lambda}+\delta\in\mathcal S\) and
\(h\leq h_1,h_2\). The minimum \(g\) is used only as a continuous
comparison function; it is never asserted to belong to \(\mathcal H\).

Ordered by decreasing pointwise value, \(\mathcal S\) is thus a net
in \(\mathcal H\). Its infimum is at least \(f\).
Conversely, at any \(x\in K\) and any real \(b>f(x)\), choose \(j\)
with \(f_j(x)<b\), then a positive rational \(\varepsilon\) and an
index \(\lambda\) with \(h_{j,\lambda}(x)+\varepsilon<b\).
The resulting member of \(\mathcal S\) is less than \(b\) at \(x\).
This proves \(\inf\mathcal S=f\), including at points where
\(f(x)=-\infty\). Excluding the identically infinite negative function
gives the conclusion.
\end{proof}

Applied to continuous MP model psh functions, the lemma proves exactly
the decreasing closure needed below from the definition of the MP cone.
Rational translation preserves that model cone. Finite maxima of general
MP psh functions are psh by the maximum property following
\cite[Definition~4.3.1]{MP24}. Bounded MP psh functions have finite energy:
monotonicity and the normalization \(E(a)=a\) bound their energy between
any constant lower and upper bounds.

\begin{lemma}\label{lem:p143k-dxz-truncation}
Let \(\Gamma\in\PSH^{\NA}_{\mathrm{DXZ}}(X,\omega)\), and let
\(\mathbf a\) denote the constant DXZ object of value \(a\in\mathbf R\).
For each positive integer \(j\) put
\(\Gamma_j=\Gamma\vee\mathbf{-j}\).
Then \(\Gamma_j\) has finite energy and
\begin{equation}\label{eq:p143k-divisorial-truncation}
 F_{\Gamma_j}(q)=\max\{F_\Gamma(q),-j\}\quad(q\in D),\qquad
 c_j=\max\{c_\Gamma,-j\}.
\end{equation}
The objects \(\Gamma_j\) decrease with \(j\).
\end{lemma}

\begin{proof}
The maximum, its order compatibility, the endpoint formula, and the
nonzero divisorial formula are
\cite[Definition~4.7 and Corollary~4.8]{Xia23}; the value at \(o\)
is precisely the endpoint, so the displayed formula holds on all of \(D\).
It remains to verify finite energy, rather than infer it from a merely
divisorial lower bound.

In every perturbation class \(\omega+\eta\), the constant object
\(\mathbf{-j}\) has level zero for \(\tau<-j\).
Since \(\Gamma_j\geq\mathbf{-j}\) and normalized model levels are at
most zero, the corresponding levels of \(\Gamma_j\) equal zero there.
Its \(-\infty\)-level volume in that component is therefore
\(\int_X(\omega+\eta)^n\). Passing to a decreasing K\"ahler perturbation
gives \(\operatorname{vol}(\Gamma_j)=V>0\).
By the positive-volume recovery theorem
\cite[Lemma~3.10]{Xia23}, \(\Gamma_j\) comes from an unperturbed
\(\mathcal I\)-model test curve \(\gamma_j\), with
\(\gamma_{j,\tau}=0\) for \(\tau<-j\).
The latter equality also follows directly from the recovery formula,
as the decreasing perturbation levels are all zero.

Using the energy definition \cite[formula~(43)]{DXZ23},
\[
 I^{\NA}_\omega(\Gamma_j)
  =c_j+\frac1V\int_{-j}^{c_j}
        \left(\int_X(\omega+\ddc\gamma_{j,\tau})^n-V\right)\rd\tau.
\]
The integral is over a finite interval; its integrand lies in \([-V,0]\).
Thus \(-j\leq I^{\NA}_\omega(\Gamma_j)\leq c_j\), as required.
This argument does not assign an unperturbed component to an arbitrary
zero-volume inverse-system object.
\end{proof}

We record the endpoint normalization in the finite-energy comparison.
For a normalized \(\mathcal I\)-maximal test curve \(\gamma\) with
endpoint \(c\), its inverse Legendre ray satisfies
\[
 U_t=\sup_{\tau<c}(\gamma_\tau+t\tau),\qquad
 \sup_X U_t=\sup_{\tau<c}(\sup_X\gamma_\tau+t\tau)=tc\quad(t>0).
\]
Each genuine model level has supremum zero. The slope identity in
\cite[Lemma~5.3.1]{MP24} therefore gives \(\phi(o)=c\) for the
corresponding MP potential. In translating conventions, both MP's
ordinary-norm potential and its logarithmic time are multiplied by two
to give the squared-norm convention above; their quotient is unchanged.
This proves the asserted zero-valuation agreement as well as the usual
nonzero divisorial agreement.

\begin{proof}[Proof of \cref{thm:p143k-full-comparison}]
First take any DXZ object \(\Gamma\). The finite-energy correspondence
assigns an MP psh function \(\phi_j\) to each truncation in
\cref{lem:p143k-dxz-truncation}. On \(D\) its values are
\(\max\{F_\Gamma,-j\}\). The MP divisorial comparison theorem gives
\(\phi_{j+1}\leq\phi_j\) everywhere. At the zero valuation,
\(\phi_j(o)=c_j\downarrow c_\Gamma\in\mathbf R\).
Hence \(\phi=\inf_j\phi_j\) is not identically \(-\infty\), and
\cref{lem:p143k-decreasing-completion} makes it MP psh.
At every \(q\in D\) its value is \(F_\Gamma(q)\).
By \cref{lem:p143k-comparison}, \(\phi=F_\Gamma\) everywhere.
This proves that the raw transform takes the full DXZ cone into the
MP cone, not just that it has a psh extension from a dense set.

Conversely, let \(\phi\) be MP psh and set
\(\phi_j=\max\{\phi,-j\}\). These are bounded psh functions, since
upper semicontinuity on the compact spectrum supplies a finite upper
bound for \(\phi\). Thus they have finite energy.
Let \(\Gamma_j\) be their inverse images under the finite-energy
correspondence. Divisorial order reflection gives
\(\Gamma_{j+1}\leq\Gamma_j\).
Their endpoints are
\[
 c_j=\phi_j(o)=\max\{\phi(o),-j\}\geq\phi(o)>-\infty.
\]
Here finiteness at \(o\in D\) is
\cite[Theorem~4.3.3]{MP24}.
The decreasing DXZ infimum exists by
\cite[Definition~4.9]{Xia23}; call it \(\Gamma\).
By \cite[Lemma~4.10]{Xia23},
\(F_\Gamma(q)=\inf_j\phi_j(q)=\phi(q)\) for every \(q\in D\),
including the endpoint formula at \(o\).
The first part gives an MP psh function \(F_\Gamma\), and
\cref{lem:p143k-comparison} identifies it with \(\phi\) on all of
\(X^\beth\). This proves surjectivity.

The DXZ divisorial order criterion makes the map injective and
order-reflecting. It is order-preserving as well: if
\(\Gamma\leq\Lambda\), their divisorial values are ordered, so the
MP comparison theorem applied to the two resulting psh functions gives
\(F_\Gamma\leq F_\Lambda\) everywhere.

On the finite-energy domain the map agrees with the existing comparison
by its divisorial values. If \(F_\Gamma\) has finite MP energy, its
finite-energy DXZ inverse image has the same divisorial values as
\(\Gamma\); injectivity makes these the same object. Thus finite-energy
domains correspond in both directions. Their energies agree by the
known finite-energy comparison; both extended energies are
\(-\infty\) on the complement. This finishes the proof.
\end{proof}

\begin{example}\label{ex:p143k-truncation-order}
For a constant DXZ object \(\mathbf a\), its raw function is the
constant \(a\). The construction uses constants \(\max\{a,-j\}\),
which are eventually exactly \(a\), and has energy \(a\).
More generally, for any \(\Gamma\), the full comparison identifies its
bounded truncation with \(\max\{F_\Gamma,-j\}\) on every semivaluation:
these are two MP psh functions with the same divisorial values.
At a point where \(F_\Gamma=-\infty\), the values of these truncations
are \(-j\) and tend to \(-\infty\). No finite-energy bound uniform in
\(j\) is required, which is why the argument reaches the full cone.
\end{example}
